\documentclass[11pt, a4paper]{article}

% Gemeinsame Style-Definitionen (TEXINPUTS muss templates/ enthalten — siehe scripts/build_tex.py)
% bfw-style.tex — gemeinsame Style-Definitionen für alle Handouts/Aufgabenhefte
% Voraussetzung: Kompilation mit xelatex (wegen fontspec)

\usepackage[ngerman]{babel}
\usepackage{geometry}
\geometry{a4paper, top=2.2cm, bottom=2.2cm, left=2.3cm, right=2.3cm}

\usepackage{fontspec}
% Fallback-Kette: Source Sans Pro -> Calibri -> Segoe UI -> Latin Modern Sans
% (Latin Modern ist immer mit MiKTeX vorhanden)
\IfFontExistsTF{Source Sans Pro}{%
  \setmainfont{Source Sans Pro}[Ligatures=TeX]%
}{\IfFontExistsTF{Calibri}{%
  \setmainfont{Calibri}[Ligatures=TeX]%
}{\IfFontExistsTF{Segoe UI}{%
  \setmainfont{Segoe UI}[Ligatures=TeX]%
}{%
  \setmainfont{Latin Modern Sans}[Ligatures=TeX]%
}}}
\IfFontExistsTF{Source Code Pro}{%
  \setmonofont{Source Code Pro}[Scale=0.92]%
}{\IfFontExistsTF{Consolas}{%
  \setmonofont{Consolas}[Scale=0.92]%
}{%
  \setmonofont{Latin Modern Mono}[Scale=0.92]%
}}

\usepackage{xcolor}
% Farbpalette: hell, freundlich, modern
\definecolor{bfwPrimary}{HTML}{0EA5A4}    % Petrol/Türkis - Primär
\definecolor{bfwPrimaryDark}{HTML}{0F766E} % dunkler Petrol für Headings
\definecolor{bfwAccent}{HTML}{F59E0B}     % Amber - Akzent
\definecolor{bfwSuccess}{HTML}{10B981}    % Grün - Erfolg / Lösung
\definecolor{bfwWarn}{HTML}{F97316}       % Orange - Achtung
\definecolor{bfwInk}{HTML}{0F172A}        % fast Schwarz
\definecolor{bfwInkSoft}{HTML}{475569}    % gedämpftes Grau
\definecolor{bfwBgSoft}{HTML}{F8FAFC}     % off-white
\definecolor{bfwBgTeal}{HTML}{ECFEFF}     % helles Türkis
\definecolor{bfwBgAmber}{HTML}{FEF3C7}    % helles Gelb
\definecolor{bfwBgRose}{HTML}{FFE4E6}     % helles Rosé für Eselsbrücken

\usepackage[colorlinks=true, linkcolor=bfwPrimaryDark, urlcolor=bfwPrimaryDark]{hyperref}
\usepackage{enumitem}
\setlist[itemize]{leftmargin=*, itemsep=2pt, topsep=2pt}
\setlist[enumerate]{leftmargin=*, itemsep=2pt, topsep=2pt}

\usepackage[most]{tcolorbox}
\usepackage{booktabs}
\usepackage{tikz}
\usetikzlibrary{decorations.pathmorphing, positioning, shapes.geometric, arrows.meta}

% Merkkästchen — wichtige Definition / Kernpunkt
\newtcolorbox{merksatz}[1][]{%
  enhanced, breakable,
  colback=bfwBgTeal,
  colframe=bfwPrimary,
  boxrule=0.6pt,
  arc=4pt,
  left=10pt, right=10pt, top=8pt, bottom=8pt,
  fonttitle=\bfseries\color{bfwPrimaryDark},
  title={\faLightbulb\ Merksatz},
  #1
}

% Eselsbrücke — Mnemoniks
\newtcolorbox{eselsbruecke}[1][]{%
  enhanced, breakable,
  colback=bfwBgRose,
  colframe=bfwAccent,
  boxrule=0.6pt,
  arc=4pt,
  left=10pt, right=10pt, top=8pt, bottom=8pt,
  fonttitle=\bfseries\color{bfwWarn},
  title={Eselsbrücke},
  #1
}

% Beispiel-Box
\newtcolorbox{beispiel}[1][]{%
  enhanced, breakable,
  colback=bfwBgSoft,
  colframe=bfwInkSoft,
  boxrule=0.4pt,
  arc=3pt,
  left=10pt, right=10pt, top=8pt, bottom=8pt,
  fonttitle=\bfseries\color{bfwInk},
  title={Beispiel},
  #1
}

% Lösungs-Box (für Aufgabenhefte)
\newtcolorbox{loesung}[1][]{%
  enhanced, breakable,
  colback=bfwBgAmber!40,
  colframe=bfwSuccess,
  boxrule=0.6pt,
  arc=3pt,
  left=10pt, right=10pt, top=8pt, bottom=8pt,
  fonttitle=\bfseries\color{bfwSuccess!70!black},
  title={Lösung},
  #1
}

% Glossar-Eintrag
\newcommand{\glossareintrag}[2]{%
  \par\noindent\textbf{\color{bfwPrimaryDark}#1} \enspace #2\par\smallskip
}

% Section-Stil — etwas farbiger als Standard
\usepackage{titlesec}
\titleformat{\section}
  {\Large\bfseries\color{bfwPrimaryDark}}
  {\thesection}{0.6em}{}
\titleformat{\subsection}
  {\large\bfseries\color{bfwPrimary}}
  {\thesubsection}{0.5em}{}
\titleformat{\subsubsection}
  {\normalsize\bfseries\color{bfwInk}}
  {\thesubsubsection}{0.4em}{}

% TikZ-Stil "skizzenhaft" — etwas weicher als Rough.js, aber stabil
\tikzset{
  roughbox/.style = {
    draw=bfwPrimaryDark, thick, fill=bfwBgTeal,
    rounded corners=4pt, inner sep=6pt
  },
  roughaccent/.style = {
    draw=bfwAccent, thick, fill=bfwBgAmber,
    rounded corners=4pt, inner sep=6pt
  },
  roughline/.style = {
    draw=bfwInkSoft, thick, -{Stealth[length=2.5mm]}
  }
}

% Bullet-Symbole (bewusst kein FontAwesome — vermeidet Font-Installations-Probleme)
\newcommand{\faLightbulb}{$\bullet$}
\newcommand{\faBookmark}{$\bullet$}

% Header/Footer
\usepackage{fancyhdr}
\pagestyle{fancy}
\fancyhf{}
\renewcommand{\headrulewidth}{0pt}
\fancyfoot[L]{\small\color{bfwInkSoft}\thepage}
\fancyfoot[R]{\small\color{bfwInkSoft} BFW M-064 Beschaffung im Büromanagement}


\title{\vspace*{-1.5cm}\Huge\bfseries\color{bfwPrimaryDark}Tag~6: Terminplanung \& Terminmanagement\\[0.4em]
  \large\color{bfwInkSoft}\textnormal{Termine im Griff --- von der Wiedervorlage bis zum geteilten Kalender}}
\author{\textsf{Büroprozesse gestalten und Arbeitsvorgänge organisieren \textperiodcentered{} M-067}\\
\textsf{\small\copyright\ Can Siebert\,\textperiodcentered\,2026}}
\date{06.07.2026}

\begin{document}
\maketitle
\thispagestyle{empty}

\vspace{1em}
\begin{merksatz}[title={\bfseries Worum geht es heute?}]
Kein Arbeitstag im Büro kommt ohne Termine aus --- Besprechungen, Liefertermine,
Steuerzahlungen, Urlaub. Wer Termine plant, überwacht und koordiniert, hält den Betrieb
am Laufen; wer sie vergisst, riskiert Mahngebühren, Vertrauensverlust und Chaos. In
diesem Handout lernen Sie, was einen Termin ausmacht, welche Terminarten es gibt, wie
Terminüberwachung und Wiedervorlage funktionieren, wie Urlaub rechtssicher geplant wird
und mit welchen Hilfsmitteln --- vom Papierkalender bis zu Outlook, geteilten Kalendern
und Doodle --- Sie Termine effizient managen. Das ist die direkte Grundlage für Tag~7
(Sitzungen \& Besprechungen) und für die IHK-Abschlussprüfung.
\end{merksatz}

\tableofcontents
\vspace{1em}
\hrule
\vspace{1em}

\section{Was ist ein Termin?}

Ein \textbf{Termin} ist ein genau definierter Zeitpunkt, an dem etwas zu erledigen ist
oder ein Ereignis stattfindet. Er besteht aus den beiden Angaben \emph{Kalenderdatum}
und \emph{Uhrzeit}. Termine im betrieblichen Kontext sind z.\,B.\ Besprechungen,
Geschäftsreisen, Urlaub, Messen, Vertreterbesuche oder Zahlungs- und Liefertermine.

Die Planung der Arbeitszeit, der Anwesenheit bei Terminen sowie die Überwachung und
Einarbeitung von Terminen in den Tagesablauf sind wesentliche Teile der Büroarbeit.
Bevor ein Termin angemessen eingeplant werden kann, sollte man sich über die Details
Gedanken machen. Dabei helfen fünf \emph{W-Fragen}:

\begin{itemize}
  \item \textbf{Warum} findet der Termin statt?
  \item \textbf{Wer} nimmt an dem Termin teil?
  \item \textbf{Wann} findet der Termin statt? (Datum, Uhrzeit)
  \item \textbf{Wo} findet der Termin statt? (Ort, Raum)
  \item \textbf{Wie lange} dauert der Termin?
\end{itemize}

\begin{beispiel}
Herr Klein aus der Abteilung Vertrieb hat am 12.\,Juni um 11:00 Uhr eine Besprechung mit
seinen Mitarbeitern. --- Die Rechnung des Holzlieferanten Opitz~AG muss bis zum
30.\,September bezahlt sein. Beide Termine bestehen aus Datum und (Spätestens-)Zeitpunkt.
\end{beispiel}

\begin{eselsbruecke}
\textbf{Die fünf W: „Warum, Wer, Wann, Wo, Wie lange?”} Wer diese fünf Fragen vor jedem
Termin beantwortet, vergisst weder Raumbuchung noch Teilnehmer noch Pufferzeit. Oft wird
zur visuellen Verdeutlichung außerdem mit einem \emph{Netzplan} gearbeitet.
\end{eselsbruecke}

\section{Terminarten}

Termine lassen sich grundlegend in \textbf{feste} und \textbf{variable} Termine
unterscheiden.

\subsection{Feste (unveränderliche) Termine}

Die \emph{festen} Termine wiederholen sich periodisch. Sie werden möglichst frühzeitig
--- teilweise schon am Jahresanfang --- in den Kalender eingetragen, um späteren
Terminüberschneidungen vorzubeugen. Feste bzw.\ unveränderliche Termine sind z.\,B.:

\begin{itemize}
  \item Geburtstage
  \item Steuerzahlungen
  \item Messen und Ausstellungen
  \item Schul- und Betriebsferien
  \item Jubiläen
  \item Hauptversammlungen
\end{itemize}

\subsection{Variable (bewegliche) Termine}

Die \emph{variablen} Termine werden sofort eingetragen, nachdem sie vereinbart wurden.
Bei Bedarf muss eine Abstimmung mit den bereits festliegenden Terminen erfolgen. Zu den
beweglichen Terminen gehören beispielsweise:

\begin{itemize}
  \item Besprechungen, Sitzungen, Tagungen
  \item Geschäftsreisen
  \item private und geschäftliche Treffen (Geschäftsessen, Konzertbesuche etc.)
\end{itemize}

\medskip
Zusätzlich unterscheidet man \textbf{Einzeltermine} (einmalig) und \textbf{Serientermine}
(nach einem Muster wiederkehrend, z.\,B.\ wöchentliches Teammeeting) sowie
\textbf{interne} Termine (im Unternehmen) und \textbf{externe} Termine (mit Kunden,
Lieferanten oder Partnern).

\begin{merksatz}
\textbf{Fest = früh eintragen} (wiederkehrend, planbar), \textbf{variabel = sofort
eintragen und abstimmen} (neu vereinbart). Bei der Eintragung lohnt es sich, für jede
Terminart eine \emph{eigene Farbe} zu verwenden --- so hat man sofort optisch einen
Überblick.
\end{merksatz}

\section{Terminüberwachung}

Die wichtige Aufgabe der \textbf{Terminüberwachung} wird oft vom Sekretariat übernommen.
Gleichzeitig ist jeder selbst für die Einhaltung seiner Termine verantwortlich. Sinnvoll
ist es, jeweils am Ende eines Arbeitstages zu prüfen, welche Termine erledigt wurden und
was am nächsten Tag dringend zu tun ist. Diese Termine werden dann für den nachfolgenden
Arbeitstag eingeplant. Hierbei unterstützt der Einsatz der \emph{Wiedervorlagemappe}.

\subsection{Gesetzlich vorgeschriebene Termine}

Gesetzlich vorgeschriebene Termine sind unbedingt einzuhalten und deshalb besonders zu
kontrollieren. Beispiele sind die Zahlung von Sozialabgaben und Steuern:

\begin{center}
\begin{tabular}{p{8cm} p{5cm}}
\toprule
\textbf{Termin} & \textbf{Fälligkeit}\\
\midrule
Sozialabgaben (Arbeitnehmer- und Arbeitgeberanteil zur Sozialversicherung) & drittletzter Bankarbeitstag des Monats\\
Einbehaltene Lohnsteuer (an das Finanzamt) & bis zum 10.\ des Folgemonats\\
\bottomrule
\end{tabular}
\end{center}

\subsection{Fristgerecht --- oder nicht}

Wird ein Termin eingehalten, spricht man von einem \emph{fristgerechten} oder
\emph{fristgemäßen} Termin. Wird ein Termin \emph{nicht} fristgerecht eingehalten, kann
er entweder gar nicht eingehalten, verspätet eingehalten oder vertagt werden. Vergessene
oder übersehene Termine sind nicht nur peinlich und unhöflich, sie können auch
unangenehme Konsequenzen nach sich ziehen.

\begin{beispiel}
Herr Fernandez aus der Buchhaltung hat die zweite Mahnung des Lieferanten Opitz~AG auf
dem Tisch liegen. Der versäumte Zahlungstermin hat zur Folge, dass kein Skontoabzug mehr
möglich war; das Unternehmen muss eventuell Verzugszinsen und Mahngebühren zahlen, und
das bisher vertrauensvolle Verhältnis zum Lieferanten ist gestört.
\end{beispiel}

Ist eine Termineinhaltung einmal nicht möglich, sollte der Termin sofort --- am besten
telefonisch und bei der zuständigen Person --- abgesagt und ein neuer Termin vereinbart
werden. Gut ist es, wenn dabei sofort ein \emph{Alternativtermin} vorgeschlagen werden
kann.

\section{Urlaubsplanung}

Eine spezielle Art der Terminplanung ist die \textbf{Urlaubsplanung}. Dabei werden die
gewünschten oder genehmigten Urlaubstage aller Mitarbeiter einer Abteilung in eine
Kalenderübersicht eingetragen. Wichtig ist es, die \emph{Rahmenbedingungen} der
einzelnen Abteilungen und --- falls möglich --- persönliche Wünsche zu beachten.
Zusätzlich muss der \emph{§\,7 des Bundesurlaubsgesetzes (BUrlG)} beachtet werden:

\begin{merksatz}[title={\bfseries §\,7 BUrlG (Kernaussagen)}]
\textbf{(1)} Bei der zeitlichen Festlegung des Urlaubs sind die Urlaubswünsche des
Arbeitnehmers zu berücksichtigen --- es sei denn, dringende betriebliche Belange oder
unter sozialen Gesichtspunkten vorrangige Urlaubswünsche anderer Arbeitnehmer stehen
entgegen.\par\smallskip
\textbf{(2)} Der Urlaub ist zusammenhängend zu gewähren, es sei denn, dringende
betriebliche oder in der Person des Arbeitnehmers liegende Gründe machen eine Teilung
erforderlich. Bei einem Anspruch von mehr als zwölf Werktagen muss ein Urlaubsteil
mindestens zwölf aufeinanderfolgende Werktage umfassen.
\end{merksatz}

\begin{beispiel}
In der Abteilung Beschaffung (Leitung: Frau Beckmann) gelten u.\,a.\ diese
Rahmenbedingungen: Herr Anderson, Frau Elling und Herr Lembke vertreten sich gegenseitig;
inklusive der Abteilungsleitung müssen immer zwei Mitarbeiter anwesend sein; im Mai darf
wegen einer Produkteinführung kein Urlaub genommen werden; Herr Lembke und Frau Elling
haben schulpflichtige Kinder und möchten nur in den Schulferien Urlaub nehmen. Solche
Bedingungen müssen frühzeitig in der Kalenderübersicht abgestimmt werden.
\end{beispiel}

\section{Hilfsmittel zur Terminplanung}

Neben den \textbf{klassischen} Hilfsmitteln wie Kalender, Planer, Terminmappen oder
Stecktafeln werden in der betrieblichen Praxis vor allem \textbf{elektronische}
Terminkalendersysteme verwendet. Elektronische Kalender gibt es als eigenständige
Programme, oftmals sind sie jedoch bereits in ein E-Mail-Programm wie z.\,B.\ MS~Outlook
integriert.

\subsection{Mindestfunktionen und Vergleich}

Folgende Funktionalitäten sollte ein Programm zur elektronischen Terminplanung
mindestens enthalten: \emph{Kalenderfunktion, Terminüberschneidungsprüfung, Adressbuch,
Notizen, Erinnerungsfunktion} sowie die Möglichkeit, den Kalender mit einem
\emph{Passwort} zu schützen.

\begin{tabular}{p{6.2cm} p{6.2cm}}
\toprule
\textbf{Vorteile elektronischer Terminplaner} & \textbf{Nachteile}\\
\midrule
Kalender über Netzwerk verbunden: Termine austauschen, abgleichen, gemeinsam bearbeiten & Gefahr des Datenverlustes bei Hard- oder Softwareproblemen\\
Frei/Belegt-Markierung; automatische Suche nach gemeinsamen freien Terminen & Änderungen sind im Nachhinein nicht mehr nachvollziehbar\\
Erinnerung per optischer oder akustischer Meldung; leichtes Ändern, Verschieben, Löschen & Ohne Netzverbindung muss synchronisiert werden\\
Mobiler Zugriff per Notebook/Smartphone; Schutz durch Passwort & Anschaffungs- und Einarbeitungsaufwand\\
\bottomrule
\end{tabular}

\subsection{Microsoft Outlook}

Stellvertretend für einen elektronischen Terminkalender wird hier \textbf{Microsoft
Outlook} vorgestellt. Outlook ist Bestandteil von Microsoft Office, zählt zu den
\emph{Personal Information Managern (PIM)} und unterstützt neben einer E-Mail-Funktion
vor allem die Terminverwaltung. Für die volle Funktionalität ist der Betrieb eines
\emph{Exchange-Servers} notwendig, der sich um die zentrale Ablage und Verwaltung von
E-Mails, Terminen, Adressen und Ressourcen kümmert. Outlook definiert für seine
Kalenderfunktion vier verschiedene \emph{Terminarten}:

\begin{itemize}
  \item \textbf{Ressourcen} --- Räume, Fahrzeuge oder Technik (Laptop, Beamer), deren
    Reservierung zeitlich genau festgelegt werden muss.
  \item \textbf{Termine} --- Anlässe ohne weitere Personen und ohne Ressource
    (z.\,B.\ Anrufe, Arztbesuch).
  \item \textbf{Besprechungen} --- Einladung weiterer Personen, zusätzlich wird eine
    Ressource (z.\,B.\ Besprechungsraum) benötigt.
  \item \textbf{Ereignisse} --- Anlässe, die 24 Stunden oder länger andauern und keiner
    bestimmten Uhrzeit zugeordnet sind (z.\,B.\ Geburtstag, Feiertag, Urlaub). Outlook
    erinnert standardmäßig 18 Stunden vorher.
\end{itemize}

\subsubsection{Besprechungsanfrage in Outlook erstellen}

Oft ist es notwendig, für eine Besprechung mit mehreren Teilnehmern einen gemeinsamen
Termin zu finden. Outlook bietet dafür die \emph{Besprechungsanfrage}:

\begin{enumerate}
  \item MS~Outlook öffnen und in die Kalenderansicht wechseln.
  \item Auf „Neue Besprechung” klicken; einen Titel (Betreff) und einen Ort eingeben.
  \item Datum, Beginn und Ende eingeben (bei ganztägigem Termin Haken bei „Ganztägiges
    Ereignis”); optional einen individuellen Text einfügen.
  \item Erforderliche und optionale Teilnehmer aus dem Adressbuch auswählen.
  \item Über die Schaltfläche „Terminplanung” die Verfügbarkeit aller Teilnehmer prüfen
    und die Zeit mithilfe der grünen und roten Linie anpassen.
  \item Auf „Senden” klicken --- die Besprechungsanfrage geht an alle Teilnehmer.
\end{enumerate}

Außerdem können Teilnehmer zu \emph{Skype-} oder \emph{Teams}-Besprechungen eingeladen
werden. Über \textbf{Serientypen} lässt sich einstellen, ob ein Anlass täglich,
wöchentlich (z.\,B.\ Teammeeting), monatlich oder jährlich stattfindet. Mit
\textbf{Terminkategorien} werden Termine farblich gekennzeichnet (z.\,B.\ Fortbildung,
Kundentermin, Lieferantentermin, Private Termine, Teammeeting, Wichtiger Abgabetermin).

\subsection{Geteilte Kalender und Online-Terminfindung}

Sollen Termine mit externen Partnern ohne MS~Outlook oder im privaten Bereich geplant
werden, hilft ein einfaches, kostenloses \emph{Online-Tool zur Terminfindung} wie
\textbf{Doodle} oder der \textbf{Google Kalender}. Doodle ermöglicht mit einem
selbsterklärenden Assistenten das Erstellen von Terminabfragen; Doodle und Outlook
lassen sich miteinander verbinden (bei Outlook.com und Microsoft~365 mit einem Klick,
sonst über ICS-Feeds).

Moderne Teams arbeiten zunehmend mit \textbf{geteilten Kalendern} (z.\,B.\ in
Microsoft~365 oder Google Workspace), in denen mehrere Personen Verfügbarkeiten sehen
und Termine gemeinsam pflegen. Für die Selbstbuchung von Terminen durch Kunden oder
Bewerber kommen \textbf{Terminbuchungs-Tools} wie Calendly oder Microsoft Bookings zum
Einsatz: Andere wählen aus den freigegebenen freien Zeitfenstern selbst einen Slot ---
das spart das mühsame Hin und Her per E-Mail. Bei allen Cloud-Tools ist der
\emph{Datenschutz} zu beachten.

\section{Elektronische Besprechungen}

Besprechungen beanspruchen einen erheblichen Teil der Arbeitszeit. Deshalb planen immer
mehr Unternehmen \emph{elektronische} Formen der Besprechung, die nicht mehr die
Anwesenheit aller Betroffenen an einem Ort erfordern.

\begin{itemize}
  \item \textbf{Videokonferenz} --- Gesprächspartner an weit auseinanderliegenden Orten
    tauschen sich per Video aus. Benötigt werden spezielle Videokonferenzsysteme und ein
    schneller Internetanschluss; auch die gemeinsame Arbeit an Dateien ist möglich.
    Eingespart werden Flug-, Hotel- und Verpflegungskosten sowie Zeit --- zunächst fallen
    jedoch Anschaffungskosten für die Technik an.
  \item \textbf{Telefonkonferenz} --- ein Telefonat mit mehr als drei Teilnehmern (genau
    drei: \emph{Dreierkonferenz}). Benötigt wird spezielle Software; manche Anbieter
    stellen virtuelle Konferenzräume bereit, die nur über eine PIN „betreten” werden und
    Platz für mehrere hundert Teilnehmer bieten.
\end{itemize}

Damit ist die Brücke zu \textbf{Tag~7} geschlagen: Erst wenn ein gemeinsamer Termin
gefunden ist, kann eine Sitzung oder Besprechung inhaltlich \emph{vorbereitet und
durchgeführt} werden.

\section{Zusammenfassung}

\begin{enumerate}
  \item Ein \textbf{Termin} ist ein definierter Zeitpunkt aus \textbf{Datum und Uhrzeit}; vor dem Einplanen helfen die fünf \textbf{W-Fragen} (Warum, Wer, Wann, Wo, Wie lange).
  \item \textbf{Terminarten}: feste (früh eintragen) und variable (sofort abstimmen) Termine; zusätzlich Einzel-/Serientermine und interne/externe Termine.
  \item \textbf{Terminüberwachung} mit Wiedervorlage(-mappe) und Farbsystem; gesetzliche Termine: Sozialabgaben (drittletzter Bankarbeitstag), Lohnsteuer (bis 10.\ des Folgemonats).
  \item \textbf{Urlaubsplanung} als Sonderfall: Kalenderübersicht, Rahmenbedingungen und \textbf{§\,7 BUrlG} (Wünsche berücksichtigen, zusammenhängend gewähren).
  \item \textbf{Hilfsmittel}: klassisch (Kalender, Mappen) und elektronisch (Outlook mit vier Terminarten, Besprechungsanfrage, Serientypen, Kategorien); Vor-/Nachteile abwägen.
  \item \textbf{Moderne Tools \& elektronische Besprechungen}: geteilte Kalender, Doodle, Terminbuchungs-Tools; Video- und Telefonkonferenz sparen Kosten und Zeit.
\end{enumerate}

\newpage

\section*{Glossar}
\addcontentsline{toc}{section}{Glossar}

\glossareintrag{Termin}{Genau definierter Zeitpunkt aus Kalenderdatum und Uhrzeit, an dem etwas zu erledigen ist oder ein Ereignis stattfindet.}

\glossareintrag{W-Fragen}{Warum, Wer, Wann, Wo und Wie lange --- Leitfragen zur Klärung eines Termins vor dem Einplanen.}

\glossareintrag{Feste Termine}{Periodisch wiederkehrende, unveränderliche Termine (z.\,B.\ Steuerzahlungen, Messen), die frühzeitig eingetragen werden.}

\glossareintrag{Variable Termine}{Bewegliche Termine (z.\,B.\ Besprechungen, Geschäftsreisen), die sofort nach Vereinbarung eingetragen und abgestimmt werden.}

\glossareintrag{Serientermin}{Anlass, der nach einem Muster regelmäßig wiederkehrt --- täglich, wöchentlich, monatlich oder jährlich.}

\glossareintrag{Terminüberwachung}{Kontrolle der Einhaltung von Terminen; oft im Sekretariat, aber jeder ist auch selbst verantwortlich.}

\glossareintrag{Wiedervorlagemappe}{Ordnungshilfsmittel, das Vorgänge zum richtigen Zeitpunkt erneut zur Bearbeitung vorlegt.}

\glossareintrag{Fristgerecht}{Einhaltung eines Termins zum vereinbarten Zeitpunkt (auch: fristgemäß).}

\glossareintrag{Urlaubsplanung}{Besondere Terminplanung: Eintragung der Urlaubstage einer Abteilung in eine Kalenderübersicht unter Beachtung der Rahmenbedingungen.}

\glossareintrag{§\,7 BUrlG}{Regelung des Bundesurlaubsgesetzes: Urlaubswünsche berücksichtigen, Urlaub zusammenhängend gewähren.}

\glossareintrag{Elektronischer Terminkalender}{Software zur Terminverwaltung mit Kalenderfunktion, Überschneidungsprüfung, Adressbuch, Erinnerung und Passwortschutz.}

\glossareintrag{MS Outlook}{Personal Information Manager von Microsoft mit vier Terminarten: Ressourcen, Termine, Besprechungen, Ereignisse.}

\glossareintrag{Exchange-Server}{Zentraler Server für die gemeinsame Verwaltung von E-Mails, Terminen, Adressen und Ressourcen in Outlook.}

\glossareintrag{Besprechungsanfrage}{Outlook-Funktion zum Finden eines gemeinsamen Termins und zum Einladen von Teilnehmern.}

\glossareintrag{Doodle}{Kostenloses Online-Tool zur Terminfindung per Terminabfrage, auch mit Outlook verbindbar.}

\glossareintrag{Geteilter Kalender}{Gemeinsam genutzter Kalender (z.\,B.\ Microsoft~365, Google Workspace), in dem mehrere Personen Verfügbarkeiten sehen.}

\glossareintrag{Terminbuchungs-Tool}{Werkzeug (z.\,B.\ Calendly, Microsoft Bookings), mit dem andere selbst freie Zeitfenster buchen.}

\glossareintrag{Videokonferenz}{Besprechung per Video an verschiedenen Orten; spart Reise- und Hotelkosten sowie Zeit.}

\glossareintrag{Telefonkonferenz}{Telefonat mit mehr als drei Teilnehmern; bei genau drei spricht man von einer Dreierkonferenz.}

\end{document}
